\section*{Chapter \ref{ch:g8:solids} --- Pyramids and Cones}

\begin{solution}{exo:g8:solids:1}
Tetrahedron: $4$ \omterm{def:g5:solids:def}{faces}, $6$ edges, $4$ \omterm{def:g5:solids:def}{vertices}
($4 + 4 = 6 + 2$). Hexagonal base: $7$ \omterm{def:g5:solids:def}{faces}, $12$ edges, $7$
\omterm{def:g5:solids:def}{vertices} ($7 + 7 = 12 + 2$).
\end{solution}

\begin{solution}{exo:g8:solids:2}
$V = \frac13 \times 36 \times 10 = 120$ cm$^3$.

Rectangular base: $B = 24$ cm$^2$, so
$V = \frac13 \times 24 \times 7.5 = 60$ cm$^3$.
\end{solution}

\begin{solution}{exo:g8:solids:3}
$V = \frac13 \pi \times 25 \times 9 = 75\pi \approx 236$ cm$^3$.
\end{solution}

\begin{solution}{exo:g8:solids:4}
\omterm{def:g7:areas:prism}{Cylinder}: $\pi \times 16 \times 12 = 192\pi \approx 603$ cm$^3$.
\omterm{def:g8:solids:def}{Cone}: $\frac{192\pi}{3} = 64\pi \approx 201$ cm$^3$. Ratio: the \omterm{def:g8:solids:def}{cone}
is one third of the \omterm{def:g7:areas:prism}{cylinder}.
\end{solution}

\begin{solution}{exo:g8:solids:5}
$56 = \frac13 \times B \times 8$, so $B = \frac{56 \times 3}{8} = 21$
cm$^2$.
\end{solution}

\begin{solution}{exo:g8:solids:6}
The lateral surface unrolls into a disk sector whose \omterm{def:g3:shapes:shapes}{radius} is the
\emph{slant height} $s$ (the distance from the apex to the rim ---
that \omterm{def:g6:lines:objects}{segment} lies on the surface), not the height $h$.
\end{solution}

\begin{solution}{exo:g8:solids:7}
$V = \frac13 \times 35^2 \times 22 = \frac13 \times 1225 \times 22
= \frac{26\,950}{3} \approx 8\,983$ m$^3$ --- about $9\,000$ cubic
meters.
\end{solution}

\begin{solution}{exo:g8:solids:8}
Height: $h = \sqrt{13^2 - 5^2} = \sqrt{169 - 25} = \sqrt{144} = 12$
cm. \omterm{def:g6:measure:volume}{Volume}: $V = \frac13 \pi \times 25 \times 12 = 100\pi$ cm$^3$.
\end{solution}

\begin{solution}{exo:g8:solids:9}
Same \omterm{def:g3:shapes:shapes}{radius}, so the \omterm{def:g8:solids:def}{cone}'s \omterm{def:g6:measure:volume}{volume} is one third of the \omterm{def:g7:areas:prism}{cylinder}'s
\omterm{def:g6:measure:volume}{volume} \emph{for the same height}: the liquid reaches
$\frac{9}{3} = 3$ cm in the \omterm{def:g7:areas:prism}{cylinder}.
\end{solution}

\begin{solution}{exo:g8:solids:10}
Sand $=$ one \omterm{def:g8:solids:def}{cone}: $V = \frac13 \pi \times 4 \times 4.5 = 6\pi
\approx 18.8$ cm$^3$. The hourglass holds two such \omterm{def:g8:solids:def}{cones}, so the
sand fills exactly \omterm{def:g3:division:half}{half} of the total inner \omterm{def:g6:measure:volume}{volume}.
\end{solution}

\begin{solution}{exo:g8:solids:11}
\emph{1.} Yes: the formula $V = \frac13 Bh$ holds for any position of
the apex, as long as $h$ is the \omterm{def:g6:lines:perp}{perpendicular} distance to the base
plane. $V = \frac13 \times 36 \times 8 = 96$ cm$^3$.

\emph{2.} Base diagonal: $\sqrt{6^2 + 6^2} = 6\sqrt2$ cm. The longest
lateral edge is the hypotenuse of a \omterm{def:g6:shapes:triangles}{right triangle} with legs
$6\sqrt2$ (in the base plane) and $8$ (vertical):
\[
\sqrt{(6\sqrt2)^2 + 8^2} = \sqrt{72 + 64} = \sqrt{136} \approx 11.7
\text{ cm}.
\]
\end{solution}

\begin{solution}{pb:g8:solids:1}
\textbf{1.} $G$ belongs to the top \omterm{def:g5:solids:def}{face} $EFGH$ and to the two
side \omterm{def:g5:solids:def}{faces} $BCGF$ and $DCGH$. The three \omterm{def:g5:solids:def}{faces} \emph{not}
containing $G$ are the bottom $ABCD$, the front $ABFE$ and the
left \omterm{def:g5:solids:def}{face} $ADHE$ (they all share the \omterm{def:g5:solids:def}{vertex} $A$, the corner
opposite $G$). The three \omterm{def:g8:solids:def}{pyramids} are $ABCDG$, $ABFEG$ and
$ADHEG$: each has a square \omterm{def:g5:solids:def}{face} of the \omterm{def:g5:solids:def}{cube} as base and the far
\omterm{def:g5:solids:def}{vertex} $G$ as apex. Each is a square-based \omterm{def:g8:solids:def}{pyramid}: $5$ \omterm{def:g5:solids:def}{faces},
$8$ edges, $5$ \omterm{def:g5:solids:def}{vertices} (\cref{ex:g8:solids:count}).

\textbf{2.} The long diagonal $[AG]$ joins the only two \omterm{def:g5:solids:def}{vertices}
belonging to none of the three bases. A third-of-a-turn rotation
about the line $(AG)$ sends the \omterm{def:g5:solids:def}{cube} to itself (the three edges
leaving $A$ --- towards $B$, $D$, $E$ --- are shuffled in a
cycle, and so are the three edges arriving at $G$). It carries
the base $ABCD$ to $ADHE$, then $ADHE$ to $ABFE$, and back:
apexes fixed at $G$, bases cycling. Each \omterm{def:g8:solids:def}{pyramid} is thus carried
onto the next one: the three are identical copies.

\textbf{3.} Three identical pieces filling the \omterm{def:g5:solids:def}{cube} share its
\omterm{def:g6:measure:volume}{volume} equally:
\[
V_{\text{pyramid}} = \frac{a^3}{3} .
\]

\textbf{4.} The \omterm{def:g8:solids:def}{pyramid} $ABCDG$ has base $ABCD$, of \omterm{def:g6:measure:area}{area} $a^2$,
and its apex $G$ sits vertically above the corner $C$, at height
$CG = a$ (a vertical edge of the \omterm{def:g5:solids:def}{cube}). As in
\cref{exo:g8:solids:11}, the formula applies with the
\omterm{def:g6:lines:perp}{perpendicular} height $h = a$:
\[
V = \frac13 \times a^2 \times a = \frac{a^3}{3} ,
\]
in perfect agreement with question 3 --- the decomposition and
the formula confirm each other.

\textbf{5.} For $a = 6$:
$V = \frac{6^3}{3} = \frac{216}{3} = 72$ cm$^3$ each. The
water-pouring experiment is this decomposition made liquid: the
\omterm{def:g7:areas:prism}{prism} over the base $ABCD$ with height $a$ is the \omterm{def:g5:solids:def}{cube} itself,
and it holds exactly three pyramid-fills --- here the three
fills \omterm{def:g3:numbers:evenodd}{even} assemble geometrically into the \omterm{def:g5:solids:def}{cube}.

\textbf{6.} Joining the center $O$ to the four corners of each
\omterm{def:g5:solids:def}{face} produces six \omterm{def:g8:solids:def}{pyramids}, one per \omterm{def:g5:solids:def}{face}: square base, apex $O$.
The center is equally placed with respect to the six \omterm{def:g5:solids:def}{faces} (any
rotation or symmetry of the \omterm{def:g5:solids:def}{cube} fixes $O$ and permutes the
\omterm{def:g5:solids:def}{faces}), so the six \omterm{def:g8:solids:def}{pyramids} are identical. The height of each is
the distance from $O$ to a \omterm{def:g5:solids:def}{face}: \omterm{def:g3:division:half}{half} the side, $\frac{a}{2}$.

\textbf{7.} Six identical pieces share the \omterm{def:g5:solids:def}{cube}:
$V = \frac{a^3}{6}$ each.

\textbf{8.} Formula: base \omterm{def:g6:measure:area}{area} $a^2$, height $\frac a2$:
\[
V = \frac13 \times a^2 \times \frac{a}{2} = \frac{a^3}{6} .
\]
The two answers agree.

\textbf{9.} Both decompositions concern \omterm{def:g8:solids:def}{pyramids} of very special
proportions carved from a \omterm{def:g5:solids:def}{cube}; a slender \omterm{def:g8:solids:def}{pyramid}, a lopsided
one, or a \omterm{def:g8:solids:def}{cone} cannot be assembled this way (no three \omterm{def:g8:solids:def}{cones} fill
a \omterm{def:g7:areas:prism}{cylinder}). This is why \cref{thm:g8:solids:volume} is
\emph{admitted} at this level: the decompositions and the
water-pouring experiment make the $\frac13$ entirely believable,
and the honest general proof --- integration --- is given in the
High School volume.

\textbf{10.} $V = \frac{6^3}{6} = 36$ cm$^3$; and by the formula,
$V = \frac13 \times 36 \times 3 = 36$ cm$^3$. They agree.

\textbf{11.} In the triangle $SAB$, the midpoints of the sides
$[SA]$ and $[SB]$ are joined by a \omterm{def:g6:lines:objects}{segment} \omterm{def:g6:lines:perp}{parallel} to $(AB)$ and
of length $\frac{AB}{2} = \frac{c}{2}$
(\cref{thm:g8:midpoints:direct}). The same holds on each lateral
\omterm{def:g5:solids:def}{face}, so the four midpoints form a square of side $\frac c2$,
with sides \omterm{def:g6:lines:perp}{parallel} to the base. On the vertical line through
the apex, the cut passes at the midpoint of the height (midpoint
theorem again, in a triangle through $S$, the base center and a
base \omterm{def:g5:solids:def}{vertex}): the cutting plane sits at height $\frac h2$.

\textbf{12.} The top piece is a regular square-based \omterm{def:g8:solids:def}{pyramid} of
base side $\frac c2$ and height $\frac h2$ (its base is the cut,
its apex is $S$). Its \omterm{def:g6:measure:volume}{volume} is
\[
\frac13 \times \left(\frac c2\right)^2 \times \frac h2
= \frac13 \times \frac{c^2}{4} \times \frac h2
= \frac{1}{8} \times \frac{c^2 h}{3}
= \frac{V}{8} :
\]
one eighth of the original --- not one \omterm{def:g3:division:half}{half}.

\textbf{13.} The frustum keeps the rest:
$V - \frac V8 = \frac{7V}{8}$. The half-height cut shares the
\omterm{def:g6:measure:volume}{volume} in the ratio $1 : 7$ --- the humble-looking bottom slab
holds seven times the \omterm{def:g6:measure:volume}{volume} of the pointed top.

\textbf{14.} Below \omterm{def:g3:division:half}{half} height lies $\frac78$ of the \omterm{def:g6:measure:volume}{volume}. With
$V = \frac13 \times 35^2 \times 22 = \frac{26\,950}{3} \approx
8\,983$ m$^3$ (\cref{exo:g8:solids:7}):
\[
\frac78 \times \frac{26\,950}{3} \approx 7\,860 \approx
7\,900 \text{ m}^3
\]
to the nearest hundred --- the ``lower \omterm{def:g3:division:half}{half}'' campaign cleans
almost eight times the glass \omterm{def:g6:measure:volume}{volume} of the upper one.

\textbf{15.} Doubling every dimension multiplies the base \omterm{def:g6:measure:area}{area} by
$2^2 = 4$ and the height by $2$, hence the \omterm{def:g6:measure:volume}{volume} by
$4 \times 2 = 8 = 2^3$. Scaling by $k$ multiplies \omterm{def:g6:measure:area}{areas} by $k^2$
and one more length by $k$: \omterm{def:g6:measure:volume}{volumes} grow by the factor $k^3$. The
top piece of question 12 is a scaled copy of the whole \omterm{def:g8:solids:def}{pyramid}
with $k = \frac12$, so its \omterm{def:g6:measure:volume}{volume} is
$\left(\frac12\right)^3 = \frac18$ of the whole --- the one-line
explanation.
\end{solution}
